\begin{center}
\LARGE
User Friendly Reviews of Computer Algebra Systems
\end{center}

\begin{center}
\Large
Nicolas Robidoux

\end{center}

\noindent
Date: July 20th (Saturday) \\
Time:  08:30-08:52.5 \\

\begin{center}
\large
Abstract
\end{center}

So far, comparative reviews of computer algebra systems have, for the most part, consisted of lists
of implemented packages with a sampling of related bugs. Organizing this information into a
coherent whole must begin with an explicit model of the users and what they do with and expect
from the systems. ``Scores'' must be accompanied by a description of the attributes of the yardstick
user: her skill level, her model and use of the computer algebra system (oracle, table book,
compute engine, high level programming language, teaching aid, driver for numerical or graphics
libraries...), her willingness to hand hold and check answers (for some users, a wrong answer which
can easily be found to be wrong may be quite acceptable, unlike one which is quite difficult to
``back plug''), the size of typical problems... Some ``yardstick users'' and their attributes will be
presented. 

